\section{Materia necesaria}

\begin{materia}{Notacion, dominios y agentes}
La cadena tiene dos agentes:
\[
\mathcal{A}=\{s,r\},
\]
donde \(s\) representa al proveedor o fabricante (\emph{supplier}) y \(r\) representa al minorista o retailer.

La cantidad vendida al cliente final se denota por
\[
q\in\mathbb{R}_+,
\]
y el precio final pagado por el cliente se denota por
\[
p\in\mathbb{R}_+.
\]
La demanda inversa es una funcion \(p:\mathbb{R}_+\to\mathbb{R}_+\) dada por:
\[
p(q)=a-bq,
\qquad
a>0,\quad b>0.
\]
Equivalentemente, la demanda directa es:
\[
q(p)=a-bp,
\]
pero en esta guia se trabaja principalmente con la demanda inversa \(p(q)\), porque la variable de decision natural sera la cantidad \(q\).

Los parametros principales son:
\[
\begin{array}{c|l}
\toprule
\text{Parametro} & \text{Significado} \\
\midrule
c\in\mathbb{R}_+ & \text{costo unitario de produccion del proveedor.}\\
w\in\mathbb{R}_+ & \text{precio mayorista cobrado al minorista.}\\
k\in\mathbb{R}_+ & \text{costo operativo unitario adicional del minorista, si existe.}\\
\alpha\in(0,1) & \text{fraccion asignada al minorista en un contrato de reparto.}\\
F\in\mathbb{R}_+ & \text{tarifa fija pagada por el minorista al proveedor.}\\
\bottomrule
\end{array}
\]

La letra griega \(\Pi\) se lee ``pi'' y representa \textbf{utilidad} o beneficio economico. No es la constante \(3.1416\). En este documento:
\[
\Pi=\text{utilidad}.
\]
El subindice indica \textbf{quien recibe} esa utilidad:
\[
\begin{array}{c|l}
\toprule
\text{Simbolo} & \text{Significado} \\
\midrule
\Pi_s & \text{utilidad del proveedor/fabricante } s.\\
\Pi_r & \text{utilidad del minorista/retailer } r.\\
\Pi_{SC} & \text{utilidad total de la cadena de suministro.}\\
\bottomrule
\end{array}
\]
El subindice \(SC\) viene de \emph{supply chain}. Por definicion:
\[
\Pi_{SC}=\Pi_s+\Pi_r.
\]

El superindice indica \textbf{en que escenario} se calcula esa utilidad:
\[
\begin{array}{c|l}
\toprule
\text{Superindice} & \text{Escenario} \\
\midrule
D & \text{descentralizado: proveedor y minorista optimizan por separado.}\\
C & \text{centralizado: la cadena completa se optimiza como una sola firma.}\\
A & \text{contrato A.}\\
B & \text{contrato B.}\\
\bottomrule
\end{array}
\]
Por lo tanto, la lectura correcta es:
\[
\begin{array}{c|l}
\toprule
\text{Simbolo} & \text{Como se lee} \\
\midrule
\Pi_s^D & \text{utilidad del proveedor en el escenario descentralizado.}\\
\Pi_r^D & \text{utilidad del minorista en el escenario descentralizado.}\\
\Pi_{SC}^D & \text{utilidad total de la cadena en el escenario descentralizado.}\\
\Pi_{SC}^C & \text{utilidad total maxima de la cadena centralizada.}\\
\Pi_s^A,\Pi_r^A & \text{utilidades del proveedor y minorista bajo el contrato A.}\\
\Pi_s^B,\Pi_r^B & \text{utilidades del proveedor y minorista bajo el contrato B.}\\
\bottomrule
\end{array}
\]

No se usara \(\Pi_D\) ni \(\Pi_C\) como abreviacion, porque es ambiguo: podria referirse al proveedor, al minorista o a la cadena completa. Siempre se escribira el subindice del agente o de la cadena: \(\Pi_s\), \(\Pi_r\) o \(\Pi_{SC}\).
\end{materia}

\begin{materia}{La idea economica}
El ejercicio compara decisiones en una cadena de suministro con dos agentes:
\[
\text{Proveedor/Fabricante}
\quad\longrightarrow\quad
\text{Minorista/Retailer}
\quad\longrightarrow\quad
\text{Cliente final}.
\]

El proveedor produce con costo unitario \(c\) y vende al minorista a precio mayorista \(w\). El minorista vende al mercado a precio final \(p\), enfrentando una demanda:
\[
q=a-bp
\quad\Longleftrightarrow\quad
p=a-bq.
\]
En lo que sigue, \(a\) mide el intercepto de demanda y \(b\) mide la sensibilidad de la demanda al precio. Se asume:
\[
a>c,
\]
para que exista una cantidad eficiente positiva.

Cuando cada agente maximiza su propia utilidad aparece \textbf{doble marginalizacion}: el proveedor agrega margen sobre \(c\), y luego el minorista agrega margen sobre \(w\). Esto sube el precio final, baja la cantidad vendida y reduce la utilidad total de la cadena.
\end{materia}

\begin{materia}{Cadena descentralizada}
En la cadena descentralizada hay una decision secuencial:
\[
\text{proveedor elige } w\in\mathbb{R}_+,
\qquad
\text{minorista observa } w \text{ y elige } q\in\mathbb{R}_+.
\]
El precio final \(p\) no es una variable independiente una vez elegido \(q\), porque queda determinado por la demanda inversa:
\[
p=p(q)=a-bq.
\]

Para un \(w\) dado, el minorista maximiza:
\[
\Pi_r(q)=q(p-w)=q(a-bq-w).
\]
La expresion \(q(p-w)\) corresponde a:
\[
\text{cantidad vendida}\times\text{margen unitario del minorista}.
\]
Su condicion de primer orden es:
\[
\frac{d\Pi_r}{dq}=a-w-2bq=0
\quad\Rightarrow\quad
q(w)=\frac{a-w}{2b}.
\]

El proveedor anticipa esa respuesta y elige \(w\):
\[
\Pi_s(w)=(w-c)q(w)
=(w-c)\frac{a-w}{2b}.
\]
Aqui \((w-c)\) es el margen unitario del proveedor y \(q(w)\) es la cantidad que el minorista decidira comprarle al observar \(w\).
Derivando:
\[
w_D=\frac{a+c}{2}.
\]
Luego:
\[
q_D=\frac{a-c}{4b},
\qquad
p_D=\frac{3a+c}{4}.
\]
Utilidades:
\[
\Pi_s^D=\frac{(a-c)^2}{8b},
\qquad
\Pi_r^D=\frac{(a-c)^2}{16b},
\qquad
\Pi_{SC}^D=\frac{3(a-c)^2}{16b}.
\]
\end{materia}

\begin{materia}{Cadena centralizada}
En la cadena centralizada se supone que proveedor y minorista actuan como una sola firma integrada. La unica variable de decision relevante es:
\[
q\in\mathbb{R}_+.
\]
El precio final queda dado por \(p(q)=a-bq\). La utilidad total de la cadena no incluye \(w\), porque \(w\) es una transferencia interna entre agentes. Por eso se maximiza:
\[
\Pi_{SC}(q)=q(p-c)=q(a-bq-c).
\]
La condicion de primer orden es:
\[
\frac{d\Pi_{SC}}{dq}=a-c-2bq=0.
\]
Entonces:
\[
q_C=\frac{a-c}{2b},
\qquad
p_C=\frac{a+c}{2},
\qquad
\Pi_{SC}^C=\frac{(a-c)^2}{4b}.
\]

Comparacion clave:
\[
q_C=2q_D,
\qquad
p_C<p_D,
\qquad
\Pi_{SC}^C>\Pi_{SC}^D.
\]
\end{materia}

\begin{materia}{Contrato de reparto de utilidades con \(\alpha\)}
El contrato de reparto de ingresos/utilidades busca que el minorista elija la cantidad centralizada. Definimos el ingreso total de la cadena como:
\[
R(q)=p(q)q=(a-bq)q.
\]
El ingreso marginal es:
\[
R'(q)=a-2bq.
\]

Si el minorista conserva una fraccion \(\alpha\in(0,1)\) de los ingresos y paga \(w\) por unidad:
\[
\Pi_r(q)=\alpha R(q)-wq.
\]
El minorista elige:
\[
\alpha R'(q)-w=0
\quad\Rightarrow\quad
R'(q)=\frac{w}{\alpha}.
\]
La cadena integrada exige:
\[
R'(q)=c.
\]
Para coordinar:
\[
\frac{w}{\alpha}=c
\quad\Rightarrow\quad
\boxed{w=\alpha c}.
\]

Con este contrato:
\[
\Pi_r=\alpha\Pi_{SC}^C,
\qquad
\Pi_s=(1-\alpha)\Pi_{SC}^C.
\]
El rango viable de \(\alpha\) se obtiene imponiendo que ambos agentes ganen mas que en el caso descentralizado:
\[
\alpha\Pi_{SC}^C>\Pi_r^D,
\qquad
(1-\alpha)\Pi_{SC}^C>\Pi_s^D.
\]
\end{materia}

\begin{trampa}{Plantilla de prueba}
Si cambian numeros, no memorices resultados. Memoriza esta secuencia:
\[
\begin{aligned}
&\text{descentralizado}
\rightarrow q(w)
\rightarrow w_D
\rightarrow (q_D,p_D,\Pi_s^D,\Pi_r^D,\Pi_{SC}^D),\\
&\text{centralizado}
\rightarrow (q_C,p_C,\Pi_{SC}^C)
\rightarrow \text{comparar contratos}.
\end{aligned}
\]
\end{trampa}
